\title{Byte Latent Transformer: Patches Scale Better Than Tokens}

\begin{abstract}
We present the Byte Latent Transformer ({\textsc BLT}), a novel architecture for large language models that operates directly at the byte level. This approach is the first to reach parity with tokenization-based LLMs in performance at scale, while delivering substantial gains in inference speed and model robustness. The {\textsc BLT} processes input as dynamically determined patches of bytes, using these patches as its core computational elements. Patch boundaries are dictated by the entropy of the subsequent byte, thus allocating more modeling capacity and computation where the information complexity is higher. We provide the inaugural study of byte-level model scaling controlled by floating-point operations, training configurations of up to 8B parameters and processing 4T bytes of data. Our experiments confirm that it is possible to effectively scale models trained directly on bytes, without relying on pre-defined token vocabularies. Efficiency in both training and inference is enhanced by adaptively grouping bytes into longer patches in predictable regions, accompanied by qualitative advances in complex reasoning and improved generalization to rare patterns. In sum, {\textsc BLT} exhibits a marked scaling advantage over tokenization-centric architectures at constant inference cost, by co-expanding both patch and model capacities.
\end{abstract}

\section{Introduction}
\label{section:intro}

We present the Byte Latent Transformer~(\textbf{BLT}), an architecture that eliminates the need for tokenization by operating directly on raw byte sequences. BLT is the first approach to reach parity with traditional models based on tokenization in large-scale settings, while also offering notable gains in efficiency and resilience (\S\ref{sec:robustness}). In contrast, prevailing large language models (llms) rely on a predominantly end-to-end training framework, but retain tokenization as a heuristic pre-processing phase that partitions bytes into a fixed vocabulary of tokens. This step can introduce biases in string compression and has well-documented limitations, including domain and modality dependence~\citep{dagangetting}, vulnerability to input perturbations~(\S\ref{sec:robustness}), insufficient orthographic representation~\citep{edman2024cute}, and disparities across languages~\citep{liang2023xlm,petrov2024language,limisiewicz2024myte}.

Tokenization has historically played a crucial role because training llms directly on raw bytes incurs significant computational expense, primarily owing to the extended sequence lengths involved~\citep{xue2022byt5}. Previous research has addressed this challenge by introducing more efficient forms of self-attention~\citep{el2020characterbert, clark2022canine} or adopting architectures that eliminate attention mechanisms entirely~\citep{wang2024mambabyte}~(\S\ref{section:related_work}). Nevertheless, these innovations largely benefit the training of \textit{small models}. When operating at large scale, the primary computational burden in a Transformer shifts to the expansive feed-forward network layers, which are applied to every byte, rather than to the attention modules themselves.

To optimize the use of computational resources, we introduce a dynamic, adaptive approach for clustering bytes into \textit{patches}~(\S\ref{section:patching}), along with a novel model architecture that integrates both byte-level and patch-level data. In contrast to traditional tokenization, {\textsc BLT} does not employ a predetermined patch vocabulary. Instead, it transforms arbitrary byte sequences into latent patch representations using efficient, learned encoder and decoder modules. Our results demonstrate that this strategy enables \textit{greater} computational efficiency compared to tokenization-driven models.

Tokenization-based large language models devote an equal amount of computational resources to every token, regardless of the specific difficulty or informational content of each token. This design prioritizes uniformity and simplicity at the cost of computational efficiency, as the compression methods used to determine token boundaries do not consistently align with the inherent complexity of the underlying predictions. The core principle underpinning our architecture is that computational effort should be selectively focused in regions where the prediction task demands it. For instance, predicting the conclusion of a word is generally a straightforward, low-entropy task that does not warrant the full capabilities of a large transformer, unlike the prediction of a sentence-initial word, which is often more challenging. In line with this philosophy, {\textsc BLT}'s structural design (\S\ref{section:architecture}) incorporates a hierarchy of three transformer modules: two smaller, byte-level \textit{local models}, and a more powerful, global \textit{latent transformer} (\autoref{fig:arch_high_level}). 

The decision of how bytes are grouped into patches, and thus how computational allocation is adjusted dynamically, is governed by {\textsc BLT} through the segmentation of input sequences according to the entropy in next-byte prediction. This method results in contextual groupings of bytes, each characterized by a roughly constant information density.

This work introduces the inaugural flop-constrained scaling analysis of byte-level architectures, scaling to models with up to 8B parameters and leveraging 4T training bytes. We demonstrate the feasibility of large-scale, end-to-end training directly from bytes, forgoing fixed-vocabulary tokenization. Notably, {\textsc BLT} achieves parity with Llama 3 in flop-controlled training performance,\footnote{Model computational requirements are quantified as the aggregate number of Floating Point OPerations (flops) required.} while delivering as much as a 50\% reduction in inference flops compared to baseline methods~(\S\ref{section:scaling}). 

Additionally, operating at the byte level affords notable benefits in modeling rare or long-tail distributions, leading to substantial gains. {\textsc BLT} displays heightened resilience relative to tokenizer-dependent models when faced with noisy textual data and demonstrates stronger character-level capabilities across orthography, phonological understanding, and machine translation in low-resource settings~(\S\ref{sec:robustness}). 

Moreover, with {\textsc BLT}, it becomes possible to concurrently expand both model and patch sizes within a constant inference flop envelope. Increasing patch size generally yields computation savings, which can then be reallocated to enlarge the global latent transformer, given it is invoked less frequently. Our series of scaling studies under fixed inference flops~(\autoref{fig:fixed_inference_scaling}) indicate that this approach leads to more favorable scaling behavior compared to traditional tokenization-dependent frameworks.

To conclude, our work presents several key contributions: 1) We propose {\textsc BLT}, a byte latent llm framework capable of adaptively distributing computational resources to enhance flop efficiency, 2) We demonstrate that our method can match the training flop cost of Llama 3 at the 8B scale, and also offer an adjustable trade-off that enables up to 50\% improvement in flop efficiency with only slight drops in evaluation scores, 3) The {\textsc BLT} approach introduces a novel scaling strategy for llms, permitting the expansion of model parameters while keeping inference costs constant, 4) We provide empirical evidence of {\textsc BLT} models’ superior resilience to input noise, as well as their sensitivity to sub-word features often overlooked by token-based llms. The training and inference implementations for {\textsc BLT} are available at~\url{https://github.com/facebookresearch/blt}.

\section{Patching: From Individual Bytes to Groups of Bytes}
\label{section:patching}

Dividing byte sequences into discrete \textit{patches} enables {\textsc BLT} to adaptively allocate computational resources depending on the input context. Various strategies for partitioning byte sequences into patches are illustrated in Figure~\ref{fig:patching_types}. More formally, a patching function $f_p$ decomposes a byte sequence $\pmb{x}=\{x_i,|i=1,\ldots n\}$ of length $n$ into a set of $m < n$ patches $\pmb{p}=\{p_j|j=1,\ldots,m\}$ by assigning each element $x_i$ a label from the set \{0,1\}, where a value of 1 signals the start of a new patch. The main determinant of computational expense in both token-based and patch-based systems is the number of operations performed by the core Transformer, which, in {\textsc BLT}, corresponds to the total patches required to encode data using the selected patching function. Therefore, the primary driver of computation, both during training and inference under a particular patching strategy, is the average patch length or simply the \textit{patch size}~(\S\ref{section:flops}).

We proceed by outlining three distinct patching approaches: segmenting with a constant byte count per patch~(\S\ref{section:static-patch}), segmentation at whitespace boundaries~(\S\ref{section:white-patch}), and adapting patch divisions dynamically via entropy estimates from a compact byte-level language model~(\S\ref{section:dyn-patch}). We also address the topic of incremental patching and clarify the differences between tokenization and patch-based approaches~(\S\ref{section:bpe-patch}).

\subsection{Strided Patching Every K Bytes}
\label{section:static-patch}

One of the simplest methods for organizing bytes is to divide them into patches of a constant size $k$, as implemented in MegaByte~\citep{yu2023megabyte}. This fixed-stride approach is straightforward to deploy during both training and inference, and it enables an easy way to adjust the average patch size—allowing direct management of computational cost. Nevertheless, this form of patching exhibits notable limitations. Firstly, computational resources are not adaptively distributed based on content: a transformer step $j$ may be inefficiently used, for example, in predicting uninformative tokens like whitespace in code, or may lack sufficient capacity for segments dense in information such as mathematical expressions. Secondly, this approach results in patching that is not consistent or context-aware, often causing identical sequences of bytes—for instance, the same word—to be divided in varying ways.

\subsection{Space Patching}
\label{section:white-patch}

\citet{slagle2024spacebyte} introduces a straightforward yet impactful modification to strided patching by generating new patches immediately following any space-like bytes\footnote{
    Space-like bytes refer to any byte that does not correspond to a Latin character, digit, or utf-8 continuation byte. Each patch is also required to include at least one byte that is not space-like.}, which tend to coincide with linguistically meaningful boundaries in numerous languages. Under the Space patching approach, each word is allocated a latent transformer computation step (i.e., more floating-point operations), allowing the model to consistently patch words across different sequences and concentrate computational effort on challenging prediction points, often located after space-like bytes. For instance, when completing the sentence, ``Who composed the Magic Flute?\uline{\hspace{.6em}}'', identifying the first byte of the answer is substantially more challenging compared to predicting subsequent bytes after ``M'', since the initial character drastically narrows the possible options and makes filling in ``Mozart'' relatively simple. Despite these benefits, space patching is not universally compatible across all languages and domains, and critically, lacks the capacity to adjust patch size. In the following section, we present an alternative patching strategy inspired by the observation that predicting the first byte of a word is usually the most demanding, while also enabling dynamic patch size control.

\subsection{Entropy Patching: Using Next-Byte Entropies from a Small Byte LM}
\label{section:dyn-patch}

Instead of utilizing traditional rule-based heuristics like whitespace, we adopt a data-driven methodology to pinpoint next-byte predictions with elevated uncertainty. We propose \textit{entropy patching}, a technique that determines patch boundaries by leveraging entropy measurements.

We train a small byte-level auto-regressive language model on the training data for {\textsc BLT} and compute next byte entropies under the LM distribution $p_{e}$ over the byte vocabulary $\mathcal{V}$:

\begin{equation}
H(x_i) = \sum_{v \in \mathcal{V}} p_{e}(x_i=v|\pmb{x}_{<i}) \log p_{e}(x_i=v|\pmb{x}_{<i})
\end{equation}

We evaluate two approaches for detecting patch boundaries using the entropies $H(x_i)$. The first method involves selecting points whose entropy exceeds a fixed global threshold, as shown in~\autoref{fig:patching}. The second method selects points where the entropy is significantly higher compared to the preceding value. This latter strategy can also be viewed as finding locations where the entropy interrupts an approximately monotonic decline within a patch.

\begin{align*}
    \text{Global Constraint}\hspace{1cm}H(x_t)& > \theta_g\\
    \text{Approx. Monotonic Constraint}\hspace{1cm}H(x_t)& - H(x_{t-1}) > \theta_r
\end{align*}

Patch boundaries are detected through a minimal preprocessing process that takes place while data is being loaded. This approach contrasts with~\citet{nawrot-etal-2023-efficient}, where a classifier is employed to learn entropy-based patch boundary prediction. In our experimental analysis~(\S\ref{section:experiments}), we evaluate and contrast these two techniques for separating bytes according to low and high entropy.

\subsection{The Byte-Pair Encoding (BPE) Tokenizer and Incremental Patching}
\label{section:incremental-patching}
\label{section:bpe-patch}

A wide range of contemporary llms, such as the baseline Llama 3, employ subword tokenizers like bpe~\citep{gage1994new, sennrich-etal-2016-neural}. In this context, we define ``tokens'' as byte groupings sourced from a \textit{finite} vocabulary that is fixed before training begins, whereas ``patches'' indicate dynamically formed sequences that do not rely on a predetermined vocabulary. The essential distinction lies in the fact that, for tokens, the model does not have explicit access to the raw byte-level representations.

A key advancement offered by {\textsc BLT} in comparison to models that rely on tokenization is its reconfiguration of the balance between vocabulary size and computational requirements. In conventional large language models (LLMs), an expanded vocabulary typically results in larger average tokens, reducing the total number of inference steps; however, this also necessitates a corresponding increase in the output dimensions of the model’s final projection layer. This inherent trade-off constrains tokenization-centric strategies, limiting their ability to realize notable changes in token length and inference efficiency. As an illustration, Llama 3 achieves an increase in the average token size from 3.7 to 4.4 bytes, but this improvement comes at the expense of expanding its embedding table by a factor of four relative to Llama 2.

During generation, {\textsc BLT} must determine if the current position in the byte sequence marks a patch boundary, as this decision controls whether the Latent Transformer is engaged for further computation. Importantly, this determination must be made without knowledge of future parts of the sequence that have yet to be produced. Consequently, the patching process must segment the byte sequence in a manner that does not rely on forthcoming bytes. Formally, we say that a patching scheme $f_p$ exhibits incremental patching if the following condition holds:
$$f_p(\pmb{x}_{<i}) = f_p(\pmb{x})_{<i}$$
bpe does not qualify as an incremental patching scheme, because the tokenization of a given prefix may vary based on the subsequent continuation of the sequence, thereby violating the incremental patching criterion above\footnote{Introducing a dedicated delimiter token to mark patch boundaries can convert bpe into an incremental patching scheme, but this adjustment also extends the overall byte-sequence length.}.

\section{{\textsc BLT} Architecture}
\label{section:architecture}
The {\textsc BLT} framework consists of a substantial global autoregressive language model that processes patch-based representations. Accompanying this central model are two compact local modules: one that translates byte sequences into patch representations, and another that reconstructs bytes from these patch embeddings~(\autoref{fig:arch_high_level}).

\subsection{Latent Global Transformer Model}

\textit{The Latent Global Transformer} refers to an autoregressive transformer architecture $\mathcal{G}$ comprised of $l_{\mathcal{G}}$ layers. This model takes as input a sequence of latent patch representations, $p_j$, and transforms them into output patch representations, $o_j$. Within the context of this work, patch indices are indicated with subscript $j$ and byte indices with $i$. The global model operates under a block-causal attention mechanism~\citep{dubey2024llama}, which confines attention to the current and all preceding patches within a document. Notably, this module accounts for the majority of computational cost—both in pre-training and at inference time—making the timing and frequency of its application a central lever for managing computational resources relative to the complexity of the given inputs and outputs.

\subsection{Local Encoder}
\label{section:local}

\textit{The Local Encoder Model}, represented as $\mathcal{E}$, constitutes a compact transformer-based architecture characterized by $l_{\mathcal{E}} << l_{\mathcal{G}}$ layers. Its central objective is to rapidly convert sequences of input bytes $b_i$ into high-level patch representations, $p_j$. Distinct from standard transformer designs, this model introduces cross-attention layers subsequent to each transformer layer, enabling aggregation of byte-level features into patch-level representations~(\autoref{fig:crossattn}).

Initially, each byte in the input sequence $b_i$ is projected via an embedding matrix of size $\mathbb{R}^{256 \times h_{\mathcal{E}}}$, producing the vectors $x_i$. These initial embeddings can be further enhanced using hash-based embeddings as described in~(\S\ref{sec:hash-emb}).

The representations are then processed through a stack of transformer layers interleaved with cross-attention layers~(\S\ref{sec:enc-cross-attn}), culminating in the patch-level representations $p_i$ which are subsequently handled by the global transformer, $\mathcal{G}$. Each transformer block applies a \textit{local block causal} attention mechanism, where every byte accesses a sliding window of $w_{\mathcal{E}}$ prior bytes. This attention window may extend across moving patch boundaries but is constrained by document boundaries. Further explanations of the embedding process and the structure of the cross-attention mechanism are provided in the following subsections.

\subsubsection{Encoder Hash n-gram Embeddings}
\label{sec:ngram-emb}
\label{sec:hash-emb}
An essential aspect of developing resilient and expressive representations at every position $i$ involves encoding context from earlier bytes. In {\textsc BLT}, this is accomplished by representing not only the single byte $b_i$, but also by capturing its presence within a broader byte n-gram. Specifically, for each index $i$, we initially generate corresponding byte-grams

\begin{align}
    g_{i,n}=\{b_{i-n + 1},\ldots, b_i\}
\end{align}

for every byte position $i$ and values of $n$ ranging from three up to eight.\footnote{
Byte-grams of length $n$ or larger are excluded when $i < n$. }

Subsequently, we propose hash $n$-gram embeddings, which transform all byte $n$-grams into indices of a fixed-sized embedding table $E_{n}^{hash}$ through a hashing operation, with a distinct embedding table maintained for each $n \in\{3,4,5,6,7,8\}$~\citep{bai2010learning}. 
These embeddings are summed with the corresponding byte’s embedding, after which the combined representation is normalized and forwarded as input to the local encoder. The composite embedding is computed as:

\begin{align}
e_i &= x_i + \sum_{n = 3,...,8}E_{n}^{hash}(\text{Hash}(g_{i,n})) \\
\text{where, Hash}(g_{i,n}) &= \text{RollPolyHash}(g_{i,n}) \% |E_{n}^{hash}|
\end{align}

We scale $e_i$ by dividing it by the total number of $n$-gram sizes incremented by one, employing the RollPolyHash method as specified in Appendix~\ref{appendix:rollpolyhash}. Section~\ref{section:ablations} explores how varying $n$ and the size of the embedding table for $n$-gram hash embeddings impacts trends in flop-controlled scaling laws. Besides the use of hash-based $n$-gram embeddings, we have also investigated frequency-based $n$-gram embeddings; the methodologies and findings for this line of inquiry are presented in Appendix~\ref{appendix:ngrams}.

\subsubsection{Encoder Multi-Headed Cross-Attention}
\label{sec:enc-cross-attn}
Our approach aligns closely with the input cross-attention mechanism from the Perceiver architecture~\citep{jaegle2021perceiver}, differing primarily in that, rather than using a static set of latent vectors, the latent representations in our setup are mapped to varying patch representations~(\autoref{fig:crossattn}). Additionally, attention operations are restricted to the specific bytes constituting each patch. Each patch $p_j$ is assigned a dedicated query vector, which is derived by aggregating the byte-level representations for $p_j$ via pooling and subsequently projecting them linearly through $\mathcal{E}_{C} \in \mathbb{R}^{h_{\mathcal{E}} \times (h_{\mathcal{E}} \times U_{\mathcal{E}})}$, where $U_{\mathcal{E}}$ denotes the total number of encoder cross-attention heads. More specifically, designating $f_{\text{bytes}}(p_j)$ as the byte sequence linked to patch $p_j$, the computation proceeds as follows:

\begin{align}
P_{0,j} &= \mathcal{E}_{C}(f_\text{bytes}((p_j)), f~ \text{is a pooling function} \\
P_l &= P_{l-1} + W_o\left(\text{softmax}\left(\frac{QK^T}{\sqrt{d_k}}\right)V\right) \\
\text{where } Q_j &= W_q(P_{l-1,j}), K_i = W_k(h_{l-1,i}), V_i = W_v(h_{l-1,i}) \\
h_l &= \text{Encoder-Transformer-Layer}_l(h_{l-1})
\end{align}

Here, $P \in \mathbb{R}^{n_p \times h_{\mathcal{G}}}$ denotes the set of $n_p$ patch representations that are provided as input to the global model. These patch representations are created by aggregating the byte embeddings $e_i$ associated with each patch $p_j$. The matrices $W_q$, $W_k$, $W_v$, and $W_o$ define the linear projections used to obtain queries, keys, values, and outputs, respectively. The keys and values are formed by projecting the byte-level representations $h_i$ output from the previous layer (or $e_i$ at the initial layer). A tailored masking scheme is applied such that each query $Q_j$ is restricted to attend only to those keys and values aligned with bytes within the same patch $j$. Since our architecture utilizes multi-head attention across $Q$, $K$, and $V$, and the dimension of patch representations ($h_{\mathcal{G}}$) is generally greater than $h_{\mathcal{E}}$, the matrix $P_l$ is structured as several heads of dimensionality $h_{\mathcal{E}}$ during the cross-attention step; these head outputs are then concatenated to reconstruct the dimension $h_{\mathcal{G}}$. We further apply a pre-LayerNorm transformation to the queries, keys, and values within this module and omit positional embeddings in the cross-attention process. The cross-attention block also incorporates a residual connection for improved representational flow.

\subsection{Local Decoder} 

In analogy to the local encoder, the local decoder $\mathcal{D}$ is implemented as a compact transformer architecture, consisting of $l_{\mathcal{D}} << l_{\mathcal{G}}$ layers. Its primary function is to transform the sequence of global patch embeddings, $o_j$, back into the original raw byte sequence $y_i$. The decoder predicts each byte conditioned on all previously generated bytes, utilizing the corresponding hidden states produced by the local encoder for those bytes as input. To accomplish this, it employs $l_{\mathcal{D}}$ layers that alternate between cross-attention mechanisms and transformer operations. The cross-attention layer precedes each transformer layer, serving to map global patch representations into byte-level embeddings, after which the transformer's operations refine this byte sequence further.

\subsubsection{Decoder Multi-headed Cross-Attention} 

In the decoder cross-attention, the roles of the queries and key/values are interchanged i.e. the byte-representations are now the queries, and the patch representations are now the key/values. The initial byte-representations for the cross-attention are initialized as the byte embeddings from the last encoder layer i.e. $h_{l_{\mathcal{E}}}$. The subsequent byte-representations for layer $l$, $d_{l,i}$ are computed as:

\begin{align}
D_0 &= h_{l_{\mathcal{E}}} \\
B_l &= D_{l-1} + W_o\left(\text{softmax}\left(\frac{QK^T}{\sqrt{d_k}}\right)V\right), \\
  &\text{where } Q_i = W_q(d_{l-1,i}), K_i = W_k(\mathcal{D}_C(o_j)), V_i = W_v(\mathcal{D}_C(o_j)) \\
  D_l &= \text{Decoder-Transformer-layer}_l(B_l)
\end{align}

Here, $W_k$ and $W_v$ refer to the key and value projection matrices, respectively. These matrices perform linear projections followed by a splitting operation, denoted as $\mathcal{D}_C$, applied to the final patch representations $o_j$ extracted from the global model. The query projection matrix $W_q$ acts on the byte-level representations $d_{l-1}$ obtained from the preceding decoder transformer layer (or $h_{l_{\mathcal{E}}}$ in the initial layer). The output projection, $W_o$, generates the outputs so that $B \in \mathbb{R}^{h_{\mathcal{D}} \times n_b}$, where $n_b$ indicates the number of output bytes. Subsequently, the representations at the $l$th decoder layer, $D_l$, are derived via a decoder transformer applied to the results of the cross-attention block, $B$. Similar to the cross-attention mechanism used in the local encoder, the implementation incorporates multiple attention heads, applies pre LayerNorm normalization, does not introduce positional embeddings, and utilizes a residual connection that surrounds the cross-attention component.

\section{Experimental Setup}
\label{section:experiments}
We establish a systematic experimental framework to contrast {\textsc BLT} against models based on traditional tokenization, ensuring that {\textsc BLT} does not benefit from potentially longer context windows.

\subsection{Pre-training Datasets} In this study, we utilize two primary pre-training datasets across all model scales: 1) The Llama 2 dataset~\citep{touvron2023llama}, which contains 2 trillion tokens sourced from a diverse range of publicly accessible materials, subsequently subjected to extensive cleaning and filtering procedures to enhance data quality; and 2) {\textsc BLT}-1T: a newly assembled dataset comprised of 1 trillion tokens drawn from multiple open sources, augmented with a subset of the pre-training data distributed by Datacomp-LM~\citep{li2024datacomp}. The Llama 2 dataset supports scaling law analyses, specifically leveraging the optimal token count determined by~\cite{dubey2024llama}, to inform architectural decisions for {\textsc BLT}. In contrast, the {\textsc BLT}-1T dataset is employed for comprehensive pre-training to facilitate downstream comparisons with Llama 3. Importantly, neither dataset incorporates information derived from Meta products or services. Additionally, all tokenizer-based baseline models employ the Llama 3 tokenizer, which features a 128K token vocabulary, as empirical results indicated that it yielded more robust baseline outcomes relative to the Llama 2 tokenizer.

\subsection{Entropy Model} For our experiments, the entropy model is designed as a byte-level language model, utilizing the identical training dataset as the full {\textsc BLT} framework. By default, our configuration consists of a transformer architecture comprising 100M parameters distributed across 14 layers, with each layer featuring a hidden size of 512, and implementing a sliding window attention mechanism with a span of 512 bytes. All other hyperparameters are aligned with those used in both our local and global transformer baselines. As detailed in \autoref{section:ablations}, we evaluated various model capacities, receptive field lengths, and alternative architectural choices. Notably, when the receptive field is sufficiently narrow, the parameters of the trained entropy model can be compactly represented using an efficient lookup table.

\subsection{Entropy Threshold and Equalizing Context Length} For models that employ entropy-based patching, we determine a patching threshold tailored to attain a specified mean \textit{patch size} across the pretraining data mix. Within {\textsc BLT}, as opposed to traditional tokenization, the \textit{patch size} is freely selectable, which meaningfully affects the context length accessible to the model. To preserve an equivalent average context length and prevent larger patch sizes from receiving disproportionate benefits, we fix the expected number of bytes per batch. Consequently, models configured with greater patch sizes operate with reduced sequence lengths. For Llama 2 data, an 8k byte context is used, while the {\textsc BLT}-1T dataset employs an increased context length of 16k bytes on average, all while keeping the batch size fixed at an average of 16M bytes.

Although the average number of samples per batch remains unchanged, dynamic patching approaches result in varying byte-to-patch ratios during data batching. To enhance efficiency, our implementation of {\textsc BLT} training aggregates fixed numbers of patches per batch, thereby eliminating the need for padding in the computationally intensive latent transformer component. Consequently, each batch consistently contains an identical patch count. Throughout training, we standardize the length of byte sequences by padding—and when necessary, truncating—to 12k bytes for Llama 2 datasets and 24k bytes for {\textsc BLT}-1T datasets. This strategy mitigates potential memory surges that could be caused by sequences with excessively large patches.

\subsection{Entropy Model Context}
\label{section:entropy-context}
Our empirical observations indicate that entropy patching tends to produce increasingly larger patches in tasks featuring structured content, such as multiple choice problems (see \autoref{fig:mmlu_patching} for an illustration on an MMLU example), where repetition is common. These outcomes stem from reduced entropy associated with the repeated sequences present in the entropy model context. Consequently, in the main experiment involving {\textsc BLT}-Entropy with a patch size of 4.5, we refresh the entropy context using new line markers and implement an approximate monotonicity constraint, as this strategy is less vulnerable to "entropy drift" caused by fluctuations in context length. Importantly, this adjustment pertains solely to our entropy calculations; the approach for determining the entropy threshold remains unchanged.

\subsection{FLOPs Estimation} 
\label{section:flops}

Our approach predominantly adopts the formulas for transformer flop calculations outlined in Chinchilla~\citep{hoffmann2022training}, encompassing the computational costs for feed-forward components, qkvo projections within self-attention, attention mechanism operations, and the output projection. One important distinction in our method is the treatment of the input embedding layer: we consider it as an optimized lookup operation rather than performing a dense matrix multiplication, effectively rendering it a 0-flop computation. Consistent with established practice, we also assume that the backward pass requires twice as many flops as the forward pass.

To compute flops \textit{per byte} for {\textsc BLT} models, we add up the flops for the local encoder transformer, the global latent transformer, and the local decoder transformer, together with the cross attention blocks in the encoder and the decoder:

\begin{align}
    \text{FL}_{\text{{\textsc BLT}}} &= \frac{1}{n_p} \cdot \text{Transf. FL}(h_{\mathcal{G}}, l_{\mathcal{G}}, m=n_{ctx}/n_p, V=0) \\
    &+ \text{Transf. FL}(h_{\mathcal{E}}, l_{\mathcal{E}}, m=w_{\mathcal{E}}, V=0) \\
    &+ \text{Transf. FL}(h_{\mathcal{D}}, l_{\mathcal{D}}, m=w_{\mathcal{D}}, V=256) \\
    &+ \text{Cross Attn. FL}(h_{\mathcal{E}}, l_{\mathcal{E}}, m=n_p, r=n_p/k) \cdot \frac{k}{n_p} \\
    &+ \text{Cross Attn. FL}(h_{\mathcal{D}}, l_{\mathcal{D}}, m=k, r=k/n_p)
\end{align}

Here, $n_{ctx}$ denotes the sequence length measured in bytes, while $n_p$ represents the size of each patch. The parameter $r$ is defined as the quotient between the number of queries and the number of keys/values, and $k$ specifies the ratio of patch-dimension to byte-dimension, meaning the number of split segments in the local model that are combined to create the global model representation ($k=2$ in \autoref{fig:crossattn}). $V$ stands for the vocabulary size used in the output projection, but this parameter is utilized exclusively within the local decoder. The context length, $m$, used by the attention mechanism depends on whether it is processing sequences at the byte or patch level. To account for these distinctions, we adjust the attention-related flops for each part of the model accordingly. Detailed flop calculation formulas for both Transformer-FLOPs and Cross-Attention FLOPs are outlined in Appendix~\ref{section:flops-appendix}.

\subsection{Bits-Per-Byte Estimation} Perplexity only makes sense in the context of a fixed tokenizer as it is a measure of the uncertainty for each token. When comparing byte and token-level models, following previous work~\citep{xue2022byt5, yu2023megabyte, wang2024mambabyte}, we instead report Bits-Per-Byte (BPB), a tokenizer independent version of perplexity. Specifically:

\begin{align}
    \text{BPB}(x)&= \frac{\mathcal{L}_{CE}(\pmb{x})}{\ln(2)\cdot n_{\text{bytes}}}
\end{align}

Here, the total cross-entropy loss computed over the data sample $\pmb{x}$ is divided by both the number of bytes contained in $\pmb{x}$ and a constant factor, effectively providing a normalized measure of uncertainty per byte.

\subsection{Transformer Architecture Hyperparameters} For all transformer modules within {\textsc BLT}, encompassing both the local and global variants, our configuration is primarily based on the architectural choices of Llama~3~\citep{dubey2024llama}. Specifically, we adopt the SwiGLU activation function~\citep{swiglu} in our feed-forward networks, apply rotary positional embeddings (RoPE)~\citep{RoPe} with a fixed $\theta=500000$~\citep{xiong2024effective} exclusively to the self-attention mechanism, and utilize RMSNorm \citep{RMSNorm} for the normalization of layers. For all self-attention mechanisms employing predetermined attention patterns such as \textit{block causal} or \textit{fixed-window block causal}, we implement Flash attention~\citep{dao2022flashattention}, with a fixed attention window size of 512 for those utilizing fixed-width masks. In cases where cross-attention layers require masks that dynamically depend on patch selection, we leverage Flex Attention\footnote{\url{https://pytorch.org/blog/flexattention}} to generate optimized, fused implementations, thereby greatly enhancing training efficiency.

\subsection{{\textsc BLT}-Specific Hyperparameters} To evaluate the performance of {\textsc BLT} models, we design experiments that focus on both scaling behavior and assessments on downstream tasks, selecting models with parameter counts of 400M, 1B, 2B, 4B, and 8B for comprehensive analysis. Details of each model’s architecture can be found in Appendix~Table~\ref{tab:arch}.

For initializing queries in the first cross-attention layer of the local encoder, we employ max-pooling. Each {\textsc BLT} model utilizes $500,000$ hashes produced by a single hash function, with n-gram lengths set between 3 and 8. The learning rate is fixed at $4e-4$ across all scales. We determine this rate for both token-based and {\textsc BLT} models based on a hyperparameter sweep in the range of $1e-3$ to $1e-4$ at the 400M and 1B scales, where results indicated that the same learning rate achieves optimal performance.

For experiments examining scaling with Llama-2 data, we use training batch sizes following the recommendations of~\cite{dubey2024llama}, or otherwise match the equivalent byte size. Optimization relies on the AdamW algorithm~\citep{adamw}, configuring $\beta_1 = 0.9$, $\beta_2 = 0.95$, and $\epsilon=10^{-8}$. Learning rates are scheduled with a linear warm-up across 2000 steps followed by cosine decay down to zero. We apply weight decay at a rate of 0.1, while global gradients are clipped to a maximum norm of 1.0.

\section{Scaling Trends}
\label{section:scaling}

We offer a comprehensive analysis of scaling patterns in byte-level models that can guide the continued expansion of {\textsc BLT} models. Our study addresses gaps in earlier work on byte-level models by: (a) contrasting trends within compute-optimal training configurations, (b) training corresponding 8B parameter models on substantial datasets—up to 1T tokens or 4T bytes—and assessing their performance on various downstream tasks, and (c) quantifying scaling behaviors when inference cost is held constant. Subsequent sections will explore particular benefits associated with processing byte sequences.

\subsection{Parameter Matched Compute Optimal Scaling Trends}
\label{section:parameter-matched-scaling}
We conduct experiments on the Llama 2 dataset by training multiple \textit{compute-optimal} bpe and {\textsc BLT} models, selecting four model scales that range between 1B and 8B parameters. To assess their scaling behavior, we graph the relationship between total training flops and language modeling accuracy on a representative sample of the training data mixture. The bpe models are optimized using the recommended parameter-to-dataset ratio established by Llama~3~\citep{dubey2024llama}, ensuring alignment with theoretical prescriptions for maximizing performance within a specified compute budget~\citep{hoffmann2022training}, and thus serving as a strong reference point. Corresponding to every bpe model, we also train a {\textsc BLT} counterpart on identical data, employing a Latent Transformer configuration with matched parameter count and architectural design as the respective bpe baseline.

As shown in~\autoref{fig:scaling-compute-opt} (right), {\textsc BLT} models consistently perform on par with or better than their bpe counterparts, a pattern that persists across increases in model capacity and computational cost. According to our understanding, {\textsc BLT} represents the inaugural byte-level Transformer model to demonstrate comparable scaling behavior to BPE-based approaches at compute-efficient settings. This outcome supports our hypothesis that the ideal parameter-to-compute ratio established for bpe remains relevant for {\textsc BLT}, or at minimum, does not significantly differ.

Achieving competitive bpe scaling behavior necessitates advances in both model architecture and dynamic patching techniques. In ~\autoref{fig:scaling-compute-opt} (left), we present a comparison between models employing space-patching methods and Llama 3. Our approximation of SpaceByte \citep{slagle2024spacebyte} leverages {\textsc BLT} space-patching, omitting n-gram embeddings and cross-attention modules. While SpaceByte exhibits gains relative to Megabyte, it still trails behind Llama 3's performance. The right panel of \autoref{fig:scaling-compute-opt} demonstrates that the combination of architectural refinements and dynamic patching yields notable benefits. Ultimately, {\textsc BLT} models achieve performance levels comparable to advanced tokenizer-based architectures like Llama 3, particularly at larger scales.

Additionally, our results show that for models relying on tokenization, the selection of tokenizer impacts performance: specifically, models utilizing the Llama-3 tokenizer achieve better results compared to those employing the Llama-2 tokenizer when both are trained on identical datasets.

In conclusion, our {\textsc BLT} architecture positions itself between Llama 2 and 3 with respect to scaling behaviors when operating with much larger patch sizes. Whereas the bpe tokenizers of Llama 2 and 3 yield average token lengths of 3.7 and 4.4 bytes respectively, {\textsc BLT} demonstrates comparable scaling properties using mean patch sizes of 6 and even 8 bytes. Since inference flops scale inversely with the average patch size, an 8-byte patch can result in almost a 50\% reduction in inference flops. Additionally, models adopting larger patch sizes generally exhibit enhanced performance as both model and dataset sizes are increased. Notably, {\textsc BLT} with an 8-byte patch size initially underperforms relative to the bpe Llama 2 at 1B parameters but surpasses the bpe baseline at 7B parameters. These observations imply that patch sizes of this magnitude may yield superior results at even larger scales, and suggest the potential viability of employing even greater patch sizes as both the scale of the model and computational resources expand.

\subsection{Beyond Compute Optimal Task Evaluations}
\label{section:evals}
To gain a deeper understanding of scaling behaviors, we extend training of an 8B {\textsc BLT} model beyond its compute-optimal regime, utilizing the more extensive and higher-quality {\textsc BLT}-1T dataset. We then evaluate the model’s capabilities across a selection of widely-used classification and generation benchmarks. For this analysis, we focus on the following tasks related to common sense reasoning, general knowledge, and code synthesis:
\paragraph{Classification tasks} These benchmarks comprise ARC-Easy (0-shot)~\citep{Eval_ARC}, Arc-Challenge (0-shot)~\citep{Eval_ARC}, HellaSwag (0-shot)~\citep{Eval_hellaswag}, PIQA (0-shot)~\citep{Eval_piqa}, and MMLU (5-shot)~\citep{hendrycks2020measuring}. Evaluation is carried out by scoring prompts and computing likelihoods for possible answer options, with average accuracy reported as the main metric.
\paragraph{Coding related generation tasks:}  For assessing code synthesis, we present pass@1 outcomes on MBPP (3-shot)~\citep{Eval_MBPP} and HumanEval (0-shot)~\citep{Eval_HumanEval}, quantifying the proficiency of LLMs in generating correct Python code snippets.

In \autoref{tab:evals}, we present a comparison among three models, each trained using the {\textsc BLT}-1T dataset: a model utilizing the Llama 3 tokenizer with bpe,\footnote{Our selection of the Llama 3 tokenizer, equipped with a 128k vocabulary, is based on its superior performance compared to the 32k vocabulary of Llama 2.} as well as two versions of the {\textsc BLT} model. These include one variant that applies a space-patching strategy ({\textsc BLT}-Space) and another that implements an entropy-based patching strategy ({\textsc BLT}-Entropy), with both integrating the approximate monotonicity constraint and resetting the entropy model's context upon encountering new lines, in accordance with the method outlined in~\autoref{section:entropy-context}. Training for all three models was conducted under an equivalent flop budget. For the {\textsc BLT}-Entropy model, we further fine-tuned the entropy threshold during inference, reducing it from 0.6 to 0.1, a modification that enhances task performance but requires an increased number of inference steps.

The {\textsc BLT}-Entropy model achieves superior results compared to the Llama 3 model on four of the seven tasks, despite being trained with an equal byte count. This enhancement is probably attributable to two main factors: (1) more effective utilization of training computation through dynamic patching, and (2) the explicit modeling of information at the byte level rather than at the token level.

Conversely, {\textsc BLT}-Space performs worse than the Llama 3 tokenizer on all tasks except one. Nonetheless, its greater average patch size of 6 bytes enables a marked decrease in inference flops. By contrast, both the bpe and entropy-patching models exhibit similar average patch sizes, around 4.5 bytes on the training dataset mix. Under an equivalent training budget, the model with the larger patch size processes 30\% more data compared to the other two models, which could lead {\textsc BLT} to deviate further from the compute-optimal regime.

\subsection{Patches Scale Better Than Tokens} {\textsc BLT} models enable concurrent scaling of both model and patch sizes without exceeding a fixed training and inference computational cost, and without requiring additional training data. This flexibility to increase patch size at will distinguishes patch-based approaches from traditional token-based models, which are constrained by efficiency limitations inherent to fixed vocabularies, as outlined in Section~\ref{section:bpe-patch}. Utilizing longer patches reduces computational requirements, thereby allowing excess compute resources to be invested in enlarging the global latent transformer, since this component is invoked less frequently.

A controlled inference scaling experiment is performed to examine whether increasing model size while reducing the number of steps on larger patch sizes leads to superior performance compared to smaller models that require more steps. Beginning with models containing 400 million and 3.6 billion parameters utilizing the Llama 2 tokenizer, we identify flop-matched counterparts using the Llama 3 tokenizer as well as {\textsc BLT}-Entropy architectures, both with mean patch sizes of 6 and 8 bytes, applied to the training datamix (detailed specifications provided in~\autoref{tab:fixed-inf-params}). For models with patch size 8, three encoder layers are employed in place of one. Each model is trained under several different training flop constraints.

As illustrated in \autoref{fig:fixed_inference_scaling}, {\textsc BLT} architectures demonstrate superior scaling behavior compared to tokenization-based models across both categories of inference flops. Specifically, BPE-based models exhibit an early advantage when training resources are limited, but {\textsc BLT} models soon outperform them once training extends slightly beyond the compute-optimal point. Consequently, allocating greater resources to the initial pretraining phase can yield more effective models under a fixed inference computational budget. This pattern is exemplified by 8B parameter models such as Llama 3.1, which was exposed to training data two orders of magnitude larger than what would be deemed optimal for its scale given compute constraints.

The point at which {\textsc BLT} surpasses token-based approaches in performance has moved nearer to the compute-optimal threshold when considering models in the higher flop class, decreasing from approximately 3x to 2.5x the optimal compute allocation. Likewise, the model employing a patch size of 8 displays a more pronounced scaling curve within this higher flop class, overtaking other variants earlier in the scaling regime. As highlighted in Section~\ref{section:parameter-matched-scaling}, increasing patch sizes tends to result in performance more comparable to BPE models at greater model sizes. This can be partly explained by the relatively reduced proportion of total computation expended on the byte-level Encoder and Decoder components, which do not scale up in complexity as rapidly as the Latent Transformer when increasing the model's overall parameter count. For instance, scaling total parameters twentyfold from 400M to 8B results in only a marginal doubling of the localized model parameters for {\textsc BLT}. This distinction is significant because patch size primarily influences the computation associated with the patch Latent Transformer, leaving the byte-level modules largely unaffected. This dynamic underlies the observation that {\textsc BLT}-Entropy with ps=8 increased from 1.6x to 1.7x the Llama 2 model size as we advanced to models of greater scale.

To summarize, our investigation into patch-length scaling reveals that the {\textsc BLT} patch-based framework can realize superior scaling behaviors when both the patch size and the overall model size are enlarged concurrently. These favorable trends are maintained and potentially enhanced at even greater model sizes.

\section{Byte Modeling Improves Robustness} We further evaluate the robustness of {\textsc BLT} in contrast to models utilizing tokens without explicit byte-level representation, and introduce a method for converting pretrained token-based models to operate at the byte level.
\label{sec:robustness}

\subsection{Character-Level Tasks}

One of the initial reasons for developing byte-level models was their inherent resistance to noise at the byte level within inputs and their capacity to recognize the substructure of tokens—a capability that traditional tokenizer-based approaches often lack. To quantitatively assess these attributes, we conduct supplementary experiments using benchmarks specifically designed to test resistance to input noise and sensitivity to character-level information across various languages, as well as numerical and phonemic content. The outcomes of these experiments are detailed in Table~\ref{tab:char_tasks}.

\paragraph{Noisy Data} To assess the robustness of {\textsc BLT} relative to tokenizer-based baselines, we generate perturbed variants of the standard classification tasks detailed in Section~\ref{section:evals}. Specifically, five distinct strategies for character-level noise are employed to alter the input text: (a) \textit{AntSpeak}: Reformats the text so that all letters appear in uppercase and are separated by spaces. (b) \textit{Drop}: Randomly eliminates 10\% of characters from the input. (c) \textit{RandomCase}: Randomly changes half of the characters to uppercase and the other half to lowercase across the text. (d) \textit{Repeat}: Selects 20\% of characters to repeat, up to four instances per character. (e) \textit{UpperCase}: Converts the entire text to uppercase characters. Each noising approach is independently applied during evaluation to either the prompt, the completion, or to both, with the mean performance reported across these setups. Table~\ref{tab:char_tasks} summarizes performance on noised HellaSwag~\citep{Eval_hellaswag}, showing that {\textsc BLT} consistently surpasses the tokenizer-based approaches in robustness—achieving an average margin of 8 points over its similarly trained counterpart, and exceeding the scores of Llama 3.1, which was trained on a much more extensive dataset.

\paragraph{Phonology - Grapheme-to-Phoneme (G2P)} We evaluate {\textsc BLT}'s proficiency in translating a string of graphemes (the letters or characters of a word) into its phonetic transcription (phonemes). Table \ref{tab:char_tasks} displays performance on the G2P task, utilizing a 5-shot framework with Phonology Bench~\citep{suvarna2024phonologybench}. Our findings indicate that {\textsc BLT} surpasses the standard Llama 3 1T tokenizer-based model on this task.

\paragraph{CUTE}
To evaluate character-level comprehension, we test {\textsc BLT} using the CUTE benchmark~\citep{edman2024cute}, which consists of a variety of tasks grouped into three main areas: compositional understanding, orthographic similarity, and sequence manipulation skills. This suite is particularly demanding for many tokenizer-based models. While these models often have some recognition of how their tokens are spelled, they generally find it difficult to leverage this information when asked to manipulate text at the character level. As presented in Table~\ref{tab:char_tasks}, {\textsc BLT}-Entropy surpasses both BPE Llama 3 baselines on the overall CUTE benchmark by a margin greater than 25 points. Notably, our model achieves near-perfect accuracy (99.9\%) on character manipulation and spelling-focused tasks. These significant gains are noteworthy given that {\textsc BLT} was trained on just one-sixteenth of the data volume used for Llama 3.1, suggesting that BPE tokenization presents a bottleneck for acquiring robust character-level competencies. Figure \ref{fig:cute_examples} provides examples where the Llama 3 tokenizer-based model underperforms, in contrast to our {\textsc BLT} model’s success. The only areas where BPE shows an advantage are word deletion and insertion tasks. While byte-level models may find word-level manipulations less direct, the performance gap here remains modest, implying that constructing word-level understanding from characters might be more straightforward than inferring fine-grained character operations from word tokens. For all experiments, we use the evaluation protocol and original Huggingface prompts without modification. It is possible that prompt engineering could further enhance BPE models’ performance.

\paragraph{Low Resource Machine Translation} We assess the performance of {\textsc BLT} on both directions of translation across six widely studied language families and twenty-one low-resource languages that use a range of scripts, utilizing the FLORES-101 benchmark~\citep{goyal2022flores}. The SentencePiece BLEU scores are presented in Table~\ref{tab:flores}. The findings indicate that {\textsc BLT} surpasses a model leveraging the Llama 3 tokenizer, achieving an overall gain of 2 BLEU points when translating into English and a 0.5-point increase when translating from English. On well-resourced language pairs, {\textsc BLT} matches or slightly exceeds Llama 3's performance. More notably, for a variety of language pairs among less-represented language families, {\textsc BLT} demonstrates superior results compared to Llama 3, highlighting the advantage of byte-level modeling for robust handling of diverse and rare byte sequences.

\subsection{Training {\textsc BLT} from Llama 3}
\label{sec:distillation}
We investigate an approach in which {\textsc BLT} models benefit from pre-existing, pre-trained models based on tokenizers to enhance both the speed and quality of training convergence. This is accomplished by initializing the global transformer weights in {\textsc BLT} using those from a pre-trained Llama 3.1. Following this, the global transformer parameters are updated using a learning rate set to one-tenth of that applied to both the local encoder and decoder, with the number of optimization steps matched to Llama 3’s preferred training duration. Table~\ref{tab:distill} illustrates a comparative analysis against a baseline {\textsc BLT} model. The results clearly show that initializing {\textsc BLT} with Llama 3.1 parameters leads to marked improvements over both baseline Llama 3 and {\textsc BLT}, all trained with equivalent computational cost. Additionally, as indicated in Table~\ref{tab:evals}, {\textsc BLT} models initialized from Llama 3.1 outperform our {\textsc BLT}-Entropy variant—despite the latter being trained on a much larger corpus (1T tokens)—on the MMLU benchmark. This suggests that this distillation-based initialization can significantly decrease training flops while boosting effectiveness.

This approach can be interpreted as converting models that rely on tokenizers into models that operate without them, effectively adapting a pre-trained LLaMA 3.1 model to the {\textsc BLT} framework. For thorough evaluation, we present both the original LLaMA 3.1—trained on 15T tokens—and the {\textsc BLT} variant derived from it in Table \ref{tab:distill}. When assessed, our model demonstrates a modest decrease in performance on MMLU and HumanEval benchmarks, while exhibiting more pronounced performance losses across the remaining tasks. These outcomes indicate a need for further exploration in order to make optimal use of the pre-trained model’s capabilities, particularly by enhancing strategies for constructing training data mixtures and refining key hyperparameters.

\section{Ablations and Discussion}
\label{section:ablations}
Here, we present a series of ablation studies that provide rationale for the architectural decisions made in {\textsc BLT}, as well as the selection of patching strategies and hyperparameter settings used for training the 8B parameter {\textsc BLT} model on the {\textsc BLT}-1T dataset.

\paragraph{Entropy Model Hyper-parameters} To investigate how the size of the entropy model and the length of its context window impact scaling efficiency, we conduct experiments with byte-level entropy transformer models spanning from 1m to 100m parameters, adjusting context windows from 64 up to 512 tokens. We display bits per byte (bpb) plotted against training flops for our $400m$ and $1b$ {\textsc BLT} models, both of which are trained on the Llama-2 dataset, as illustrated in \autoref{fig:entropy_ablation}. Our results demonstrate that increasing either the entropy model size or the context window length leads to improved scaling performance, though gains taper off after surpassing 50m parameters.

\paragraph{Types of Patching}

We conduct an ablation study of the four distinct patching strategies described in Section~\ref{section:patching}, namely: 1) Strided Patching employing strides of 4 and 6, 2) Patching triggered by whitespace, 3) Patching utilizing the BPE Tokenizer aligned with the Llama 3 tokenizer, and 4) Entropy-based patching leveraging a compact byte-level language model.

Although dynamic patching shortens the apparent sequence lengths, we ensure comparability by regulating the sequence length across all patching strategies, thereby preserving equivalent context lengths. Each model is exposed to an equal expected number of bytes per sequence during both training and inference phases, eliminating potential confounds related to the capacity for processing extended contexts. The outcomes of these controlled experiments are illustrated in Figure~\ref{fig:scaling-compute-opt}. All alternative patching methods demonstrate superior performance over static patching, and notably, space patching achieves results closely matching those of dynamic entropy-based patching.

As shown in \autoref{tab:evals_ablation}, we report benchmark results for {\textsc BLT} models utilizing three approaches: tokenizer-based models, space patching, and entropy-based patching, all trained for an optimal number of steps on the Llama 2 dataset~\citep{dubey2024llama}. Space patching, while more straightforward and not requiring real-time entropy model evaluation during training, does not match the improvements achieved by entropy-based patching. Notably, the benefits demonstrated by entropy-based patching for scaling performance~(see Section~\ref{section:scaling}) are maintained when evaluated on downstream benchmark tasks.\footnote{Space patching was previously tested in runs that excluded cross-attention, but similar patterns persist when cross-attention is included.}

\paragraph{Cross-Attention} 
\autoref{tab:crossattn_ablations_1b} presents ablation results on incorporating cross-attention at different stages within the encoder and decoder of {\textsc BLT}. Regarding encoder cross-attention, we evaluate three strategies for initializing the queries: (1) employing a uniform learned embedding across all global states, (2) utilizing a hash-based embedding derived from the bytes within each patch, and (3) applying pooling to the encoder's hidden states corresponding to the patch bytes at the relevant encoder layer.

Our results indicate that integrating cross-attention within the \textit{decoder} yields the highest performance gains. Employing cross-attention in the encoder offers marginal benefits, which are noticeable only when query initialization is performed via pooling. Moreover, cross-attention proves especially advantageous for the Common-Crawl dataset, with its effects being most pronounced when utilizing larger patch sizes.

\paragraph{n-gram Hash Embeddings} 

We experiment with n-gram hash embedding vocabulary sizes of 0, 100K, 200K, and 400K, with corresponding results reported in Table~\ref{tab:ngram_ablations_1b}. Our analysis reveals that the addition of hash embeddings results in performance improvements across every domain, with particularly notable effects observed for Wikipedia and Github (showing a 0.04 bpb improvement compared to a 0.01 bpb gain after 15k steps at the 8B model scale). When adjusting the number of hashes from 500K to 300K at the 8B scale, performance changed by only about 0.001 bpb after 15k steps. These findings suggest that the use of hashes plays a crucial role in enabling {\textsc BLT} models to achieve parity with tokenizer-based approaches; however, increasing the hash size beyond 300K brings only marginal benefits. Further, we observe that the performance improvements afforded by hash embeddings are mostly additive to those provided by cross-attention, as enhancements are evident on different datasets.

\paragraph{Local Model Hyperparamaters} Table \ref{tab:local_layers} presents an ablation study on the impact of varying the number of layers within the local encoder and decoder components. Results indicate that when using hash n-gram embeddings, {\textsc BLT} achieves strong performance even with a highly minimal encoder comprising only a single layer, provided that the decoder is comparatively more complex.

\section{Related Work}
\label{section:related_work}

\textbf{Character-Level RNNs:} The task of Character Language Modeling has remained prominent since the emergence of neural architectures~\citep{sutskever2011generating,mikolov2012subword,graves2013generating}, largely due to their capacity to handle out-of-vocabulary words naturally and without reliance on back-off techniques. \cite{kim2016character} introduce a model which processes input solely at the character level via convolutional and highway layers, passing these representations into LSTM-based RNNs. Their approach achieves parity with contemporary RNN-based language modeling baselines on English, while offering notable gains on languages with rich morphology—a frequently cited strength of character-level LLMs. Similarly, \cite{kenter2018byte} utilize byte-level LSTM models for machine comprehension, achieving superior results in morphologically complex languages such as Turkish and Russian compared to word-level counterparts. In related work, \cite{zhang2015character} deploy character-level convolutional networks for classification, demonstrating outperformance relative to word-based models on selected tasks. To further enhance character-level language modeling, \citet{chung2019hierarchical} propose hierarchical LSTMs with boundary detectors, facilitating the discovery of latent textual structure. Alternatively, ByteNet~\citep{kalchbrenner2016neural} approaches sequence modeling for machine translation by applying CNN-based layers on characters, foregoing attention mechanisms entirely.

\textbf{Character-Level Transformers:} The introduction of transformers leveraging attention mechanisms~\citep{vaswani2017attention} in conjunction with subword tokenization approaches~\citep{sennrich-etal-2016-neural} yielded notable advancements in the performance of neural networks for tasks such as language modeling and standardized benchmarks. Nevertheless, the use of word and sub-word representations inherently prescribes a particular inductive bias regarding the abstraction scale at which these models function. Seeking to merge the achievements of transformer architectures with the positive early outcomes seen in character-level language modeling, \citet{al2019character} implemented extremely deep transformer networks and, aided by auxiliary loss functions, developed models based on transformers that surpassed prior character-level LSTM-based language models. Despite these gains, their models continued to lag behind word-level LLMs by a notable margin. Similarly, GPT-2~\citep{radfordgpt2} reported that byte-level language models failed to match the performance of their word-level counterparts when evaluated on large-scale datasets such as the 1 billion word benchmark.

Although \cite{Choe2019BridgingTG} showed that transformer-based byte-level LLMs could surpass their subword-level counterparts with similar numbers of parameters, their approach required substantially more computational resources and significantly increased training time. In a parallel effort, \cite{el2020characterbert} introduced CharFormer, a BERT-based model that constructs word representations via convolutional operations on character embeddings and demonstrated gains in the medical domain. However, this method also entailed a considerable increase in computational cost. The CANINE model developed by \cite{clark2022canine}, an encoder-only architecture with 150M parameters processing character sequences directly, incorporates a deep transformer architecture analogous to our global model, alongside a local transformer and strided convolutions that serve to downsample input characters. This configuration allowed CANINE to outperform the equivalent token-based encoder-only model (mBERT) in downstream multilingual evaluations. In another line of work, ByT5~\citep{xue2022byt5} investigated byte-level encoder-decoder architectures that omit patching operations. Their model exhibited heightened robustness to noisy inputs and matched the performance of tokenizer-dependent models while utilizing only a quarter of the training data. However, the absence of patching imposed the need for computation-intensive attention calculations across all bytes, leading to significant increases in computational overhead. Representing sequences at the byte level—as opposed to subwords—inherently leads to much longer input sequences, posing notable efficiency challenges when scaling such models. The introduction of the Mamba Architecture~\citep{gu2023mamba}, designed to maintain a fixed-size memory even with very long input contexts, recently enabled \cite{wang2024mambabyte} to train a patch-free byte-level Mamba model. This architecture outperforms byte-level transformer models in bits-per-byte under flop-controlled conditions at the 350M parameter scale, across a range of datasets.

\textbf{Patching-based approaches:} Leveraging patching techniques has shown promise in reducing the high computational overhead typically associated with byte-level LLMs, without necessarily sacrificing model effectiveness. Early investigations demonstrated that patching can be beneficial, particularly for models with smaller architectures and datasets. For instance, \cite{nawrot-etal-2022-hierarchical} propose an hourglass transformer architecture that applies static patching through downsampling and upsampling, surpassing alternative byte-level baselines at a 150M parameter scale. Subsequently, \cite{nawrot-etal-2023-efficient} extend this work by introducing dynamic patching methodologies, which encompass an end-to-end trainable boundary-predictor, boundary-predictors trained with the guidance of specific tokenizers, and an entropy-driven patching approach akin to {\textsc BLT}. These innovations allow their models to achieve superior results compared to standard transformers of that era on language modeling benchmarks, specifically at the 40M parameter range on datasets of roughly 400M tokens. In another direction, \cite{lester2024training} explore sequence compression via arithmetic coding, attaining higher compression ratios than traditional BPE. Utilizing an equal-info windows strategy, their method surpasses standard byte-level baselines in language modeling, but remains less effective than subword-based baselines.

Our research builds upon and is most akin to MegaByte~\citep{yu2023megabyte}, a causal LLM featuring a decoder-only architecture that employs fixed, static patching along with representation concatenation to convert bytes into patches, and leverages a local model on the decoder side for reconstructing patches back into bytes. Their findings indicate that, at the 1B parameter scale and over a dataset of roughly 400B bytes, MegaByte achieves parity with models reliant on tokenizers. In all experimental analyses, we systematically include MegaByte as a baseline and observe that static patching underperforms compared to the latest state-of-the-art tokenizer-based models when training compute is held constant. We show that {\textsc BLT} effectively narrows this performance gap. Similarly, \cite{slagle2024spacebyte} arrive at comparable conclusions regarding MegaByte and propose enhancements such as performing patching not only on fixed intervals but also at whitespaces and analogous space-like bytes, along with introducing a local encoder component. Their evaluation reveals better performance over tokenizer-based transformer models within specific domains (e.g., Github and arXiv) at the 1B parameter scale under constrained compute budgets. We also provide empirical results for this extended approach, demonstrating that additional architectural advancements are necessary for byte-level models to achieve competitive scaling and match the strong results seen in current state-of-the-art token-based systems like Llama 3.

\section{Limitations and Future Work}

For this study, our architectural design decisions follow the model training step count identified as optimal for Llama~3~\citep{dubey2024llama}. Nevertheless, it is important to note that these scaling laws were originally established for transformers operating at the BPE level, and as such, their application to {\textsc BLT} may result in less-than-ideal ratios between data volume and model parameters. Deriving specific scaling laws tailored to {\textsc BLT}—which could reveal even better scaling characteristics for this architecture—is left to future research. Moreover, our experiments primarily operate at a scale of up to 1B parameters. As we expand to model sizes of 8B parameters or larger, the set of optimal architectural choices may shift, potentially yielding further performance enhancements at larger scales.

Current transformer codebases and libraries are optimized for maximal efficiency when working with tokenizer-based transformer models. Although we conduct experiments matched in theoretical FLOPs and incorporate some optimized components (like FlexAttention) to accommodate layers that depart from standard transformer designs, our implementations might still lag behind tokenizer-based approaches with respect to wall-clock performance and could see improvements through additional optimization efforts.

Whereas {\textsc BLT} employs a separately optimized entropy model for the patching component, an intriguing avenue for future exploration involves integrating patching model training into a fully end-to-end pipeline. As discussed in Section \ref{sec:distillation}, we conduct preliminary studies that suggest positive outcomes in adapting "tokenizer-based" models like Llama 3—which have been trained on datasets exceeding 10T tokens—to byte-level formats by reusing and freezing the weights of their global transformer modules. Progressing further along this line of inquiry could reveal techniques that preserve the strengths conferred by bytefying, while potentially surpassing the capabilities of the original tokenizer-based systems, all without necessitating training from the ground up.

\section{Conclusion}
\label{section:conclusion}

This work introduces the Byte Latent Transformer (\textbf{BLT}), an innovative model architecture that moves beyond the traditional reliance on static vocabularies for tokenization in large language models. Through a trainable and adaptive approach to combining bytes into larger “patches,” {\textsc BLT} enables the allocation of computational effort in proportion to the underlying complexity of the data. This approach results in marked gains in computational efficiency and the ability to generalize. Comprehensive scaling experiments reveal that {\textsc BLT} models can achieve parity with established tokenization-based counterparts such as Llama 3 at scales reaching 8B parameters and 4T bytes. Notably, they offer up to a 50\% reduction in inference flops with only slight trade-offs in conventional evaluation benchmarks. Moreover, {\textsc BLT} introduces the potential for simultaneous scaling of both model size and patch granularity without exceeding a predetermined inference compute budget—an advantage particularly relevant under real-world hardware and cost constraints. By directly processing raw byte sequences, {\textsc BLT} further enhances the model’s resilience to rare or unexpected inputs and supports more granular learning at the sub-word level. Collectively, these findings identify {\textsc BLT} as a strong alternative to established tokenization strategies, delivering a flexible and efficient foundation for robust next-generation language models.

\end{document}